% !TeX program = lualatex
% !TeX encoding = utf8
% !TeX root = ../Thermod.tex

\setcounter{chapter}{8}

%=========================================================
\Opensolutionfile{answer}[\currfilebase/\currfilebase-Answers]
\Writetofile{answer}{\protect\section*{\nameref*{\currfilebase}}}
\chapter{Поверхневий натяг}\label{\currfilebase}
\graphicspath{{\currfilebase/pictures}}
\makeatletter
\def\input@path{{\currfilebase/}}
\makeatother

%% --------------------------------------------------------
\section{Вплив поверхні на термодинаміку речовин}
%% --------------------------------------------------------

До теперішнього моменту ми нехтували впливом поверхні на термодинаміку речовин. Між тим, очевидно, що такий вплив може бути суттєвим при маленьких розмірах об’єкту, маленьких хочу б в одному вимірі. До таких об’єктів слід віднести, наприклад, порошки, тонкі плівки, капіляри тощо (рис.~\ref{pic:Pic9.1}), де енергії поверхні та енергія об'єму одного порядку $E_\Sigma \sim E_V$ або значно більше ніж в звичайних тілах. Про те, що вплив поверхні на поведінку цих об’єктів істотний, свідчать дослідні дані. Відомо, що порошки мають розвинуту поверхню, що робить їх надзвичайно активними. Особливо дрібнодисперсні порошки можуть бути вибухонебезпечні зазнаючи за рахунок великої питомої поверхні швидких хімічних перетворень (рис.~\ref{pic:Pic9.1a}) навіть якщо в об’ємному стані речовина (наприклад \ce{Ti}, \ce{Al}, \ce{Zn}) не становить небезпеки. Явище капілярного підняття примушує воду підніматися по стовбурам дерев, сформованих із найтонших капілярів, на висоту більше, ніж 10м зі швидкістю $60$~м/год (рис.~\ref{pic:Pic9.1b}). Це ж явище широко використовується у порошковій металургії, дозволяючи спікати суміші порошків речовин за температури значно нижче температури плавлення кожної з компонент (рис.~\ref{pic:Pic9.1c}).

%=========================================================
\begin{figure}[h!]\centering
%---------------------------------------------------------
\begin{subfigure}[t]{0.31\linewidth}\centering
    \includegraphics[height=4cm]{Pic9.1a}
\caption{Хімічні реакції.}
\label{pic:Pic9.1a}
\end{subfigure}
%---------------------------------------------------------
\begin{subfigure}[t]{0.31\linewidth}\centering
    \includegraphics[height=4cm]{Pic9.1b}
\caption{Капілярний ефект сприяє підняттю води в рослинах.}
\label{pic:Pic9.1b}
\end{subfigure}
%---------------------------------------------------------
\begin{subfigure}[t]{0.31\linewidth}\centering
    \includegraphics[height=4cm]{Pic9.1c}
\caption{Спікання порошків металів і сплавів.}
\label{pic:Pic9.1c}
\end{subfigure}
%---------------------------------------------------------
\caption{Визначальний вплив поверхневих явищ.}
\label{pic:Pic9.1}
\end{figure}
%=========================================================


Власне форма дрібнодисперсних об’єктів, зокрема крапельок води або кристалів льоду, що утворюють візерунки у вигляді розвинених дендритів на поверхні скла визначається силами поверхневого натягу (рис.~\ref{pic:Pic9.2}). Зрештою ключову роль сили поверхневого натягу відіграють також при фазових перетвореннях.

%=========================================================
\begin{figure}[h!]\centering
%---------------------------------------------------------
\begin{subfigure}[t]{0.45\linewidth}\centering
\includegraphics[width=\linewidth]{Pic9.2a}
\caption{}
\label{pic:Pic9.2a}
\end{subfigure}
%---------------------------------------------------------
\begin{subfigure}[t]{0.45\linewidth}\centering
\includegraphics[width=\linewidth]{Pic9.2b}
\caption{}
\end{subfigure}
%---------------------------------------------------------
\caption{Визначальний вплив поверхневих явищ на форму дрібно дисперсних об’єктів: а) крапель рідини; б) форму кристалів при фазових перетвореннях (замерзанні води).}
\label{pic:Pic9.2}
\end{figure}
%=========================================================


%До теперішнього моменту ми нехтували впливом поверхні на термодинаміку речовин. Між тим, очевидно, що такий вплив може бути суттєвим при маленьких розмірах об’єкту, або хоча б в одному напрямку. До таких об’єктів слід віднести, наприклад, порошки, тонкі плівки, капіляри тощо, де енергії поверхні та енергія об'єму одного порядку $E_\Sigma \sim E_V$, або значно більше, ніж в звичайних тілах.

Фізична причина виникнення сил поверхневого натягу полягає в тому, що збільшення поверхні, яка обмежує систему, вимагає виконання роботи. Дійсно, для збільшення площі поверхні частина молекул з об’єму має вийти на поверхню, долаючи сили молекулярного тяжіння (рис.~\ref{pic:Pic9.3}).

%---------------------------------------------------------
\begin{figure}[h!]\centering
\includegraphics[width=0.35\linewidth]{Pic9.3}
\caption{Нерівноцінність умов, в яких знаходяться молекули всередині і на поверхні.}
\label{pic:Pic9.3}
\end{figure}
%---------------------------------------------------------

 З хорошою степеню точності можна вважати, що ця робота пропорційна площі, на яку збільшилась поверхня, оскільки вона, в свою чергу, пропорційна кількості частинок, що утворюють поверхню. Тоді елементарна робота системи по утворені поверхні буде:

\begin{equation}
    \delta A = - \sigma \, d\Sigma,
    \label{eq:work_surface}
\end{equation}
де  $\sigma$ --- нова введена узагальнена сила, що відповідає узагальненій координаті $\Sigma$ --- площі. Величина ця має бути додана у вираз для диференціалів термодинамічних потенціалів, кожен раз коли слід враховувати сили поверхневого натягу. В окремому випадку для вільної енергії Гельмгольца:


\begin{equation}
    dF = -S\,dT - p\,dV + \sigma\,d\Sigma.
    \label{eq:helmholtz_surface}
\end{equation}

Звідки, за визначенням \textit{поверхневим натягом} назвемо роботу, необхідну для ізотермічного збільшення поверхні на одиницю площі при постійному об’ємі системи:

\begin{equation}
    \sigma = \left( \frac{\partial F}{\partial \Sigma} \right)_{T,V}.
    \label{eq:sigma_definition}
\end{equation}

Вимога сталості об’єму є обов’язковою, оскільки виключає інші джерела роботи. Саме робота сил поверхневого натягу визначає рівноважну форму системи при сталості її об’єму. Таким чином визначається  кулевидна форма малих крапель рідини (рис.~\ref{pic:Pic9.2a}), округла форма капель жиру в борщі, огранка  монокристалів що виростають із розплаву.

%% --------------------------------------------------------
\section{Поверхневий натяг з точки зору механіки}
%% --------------------------------------------------------


Походження назви поверхневого натягу пов’язано зі способом його визначення. Тут ідеться про принциповий спосіб, оскільки на практиці використовуються інші. Робота, яку необхідно виконати для переміщення легко ковзаючої планки довжиною $a$ на відстані $dх$, становить $F dx$. Тут $F$ сила, котру слід прикласти до планки, що утримується плівкою, натягнуту на рамку (рис.~\ref{pic:Pic9.4}). З іншого боку, за визначенням ця робота дорівнює  $\sigma a dx$. Прирівнюючи ці величини отримаємо що:

\begin{equation}
    \sigma  = \frac{F}{a}\quad \left[\frac{\text{Н}}{\text{м}}\right]\ \text{або}\ \left[\frac{\text{Дж}}{\text{м}^2}\right].
    \label{eq:force_surface}
\end{equation}

Тобто механічний сенс $\sigma$ --- це сила, віднесена до одиниці довжини поверхні. Характеристичні значення сил поверхневого натягу --- $10^{-2}$ Дж/м\textsuperscript{2}.


%---------------------------------------------------------
\begin{figure}[h!]\centering
\includegraphics[width=0.5\linewidth]{Pic9.4}
\caption{Схема принципового способу визначення коефіцієнту поверхневого натягу і його значення для деяких рідин.}
\label{pic:Pic9.4}
\end{figure}
%---------------------------------------------------------

%\begin{table}[h!]
%\centering
%\begin{tblr}{
%  colspec={|l|c|},
%  hlines,
%  vlines,
%  row{1} = {font=\bfseries},
%}
%Рідина       & Поверхневий натяг [Н/м] \\
%спирт        & 0.0223 \\
%олія         & 0.0320 \\
%мильна вода  & 0.0370 \\
%кров         & 0.0580 \\
%гліцерин     & 0.0631 \\
%вода         & 0.0728 \\
%\end{tblr}
%\end{table}


%% --------------------------------------------------------
\section{Внутрішня енергія поверхневого натягу}
%% --------------------------------------------------------

Значення внутрішньої поверхневої енергії, тобто внутрішньої енергії віднесеної до одиниці площі, знайдемо із співвідношення Гібса-Гемгольца:

\begin{equation}
    U = F - T \left( \frac{\partial F}{\partial T} \right)_{V,\Sigma}.
    \label{eq:gibbs_helmholtz}
\end{equation}

Де похідна береться, згідно \eqref{eq:helmholtz_surface} при сталих $V$ та $\Sigma$. Але згідно визначенню, вільна енергія віднесена до одиниці площі є $\sigma$. Звідки:

\begin{equation}
    U_{\Sigma}\Sigma = \sigma - T \left( \frac{\partial \sigma}{\partial T} \right)_{V,\Sigma} \Sigma,
    \label{eq:surface_energy}
\end{equation}
де $U_{\Sigma}$ --- шукана величина:

\begin{equation}
    U_{\Sigma} = \sigma - T \left( \frac{\partial \sigma}{\partial T} \right)_{V,\Sigma},
\end{equation}

Якщо плівка розшириться ізотермічно, то їй слід передати тепло, що дорівнює за першим законом термодинаміки:

\begin{equation}
    \Delta Q = \Delta U_{\Sigma} - \sigma \Delta\Sigma = (U_{\Sigma} - \sigma) \Delta\Sigma,
    \label{eq:heat_isothermal}
\end{equation}
де $\Delta \Sigma$ --- дорівнює приросту площі. Звідси з врахуванням \eqref{eq:heat_isothermal}:
\begin{equation}
    \Delta Q = q\Delta\Sigma = - T \left( \frac{\partial \sigma}{\partial T} \right)_{V,\Sigma} \Delta\Sigma.
\end{equation}

Величина $q$ --- питома теплота, яку слід передати поверхні для збільшення її на одиницю площі при постійній температурі. Ця величина додатна, оскільки похідна $\frac{\partial \sigma}{\partial T}$ від’ємна, як слідує з дослідних даних.

Величина $q$ --- питома теплота, яку слід передати поверхні для збільшення її на одиницю площі при постійній температурі. Ця величина додатна, оскільки  похідна $\frac{\partial \sigma}{\partial T}$ від’ємна, як слідує з дослідних даних (рис.~\ref{pic:Pic9.5}) Дійсно, в критичній точці, де різниця між рідиною та газом зникає, $\sigma$, очевидно, взагалі має дорівнювати нулю.

%---------------------------------------------------------
\begin{figure}[h!]\centering
\pgfplotstableread[]{
x       y1          y2          y3          y4          y5          y6          y7          y8
-20     0.00859     0.0122      0.0128      0.01494     0.01542     nan         nan         0.01440
-10     0.00650     0.0107      0.0114      0.01354     0.01429     0.01602     nan         0.01299
0       0.0045      0.00936     0.0101      0.01219     0.01301     0.01485     nan         0.01148
10      0.00273     0.00806     0.00886     0.01084     0.01186     0.01361     nan         0.01011
20      0.00118     0.00679     0.00758     0.00953     0.01066     0.01243     0.01600     0.00873
30      0           0.00557     0.00638     0.00818     0.00955     0.01126     0.01494     0.00736
40      nan         0.00444     0.00523     0.00698     0.00840     0.01011     0.01389     0.00612
50      nan         0.00337     0.00408     0.00574     0.00734     0.00895     0.01281     0.00496
60      nan         0.00240     0.00304     0.00454     0.00625     0.00787     0.01177     0.00368
70      nan         0.00146     0.00206     0.00346     0.00516     0.00680     0.01070     0.00259
80      nan         0.000755    0.00113     0.00246     0.00417     0.00581     0.00968     0.00162
90      nan         0.000155    0.000355    0.00144     0.00324     0.00481     0.00866     0.00071
100     nan         nan         nan         0.00058     0.00237     0.00386     0.00769     0.00000
110     nan         nan         nan         nan         0.00153     0.00295     0.00671     nan
120     nan         nan         nan         nan         0.00080     0.00211     0.00576     nan
130     nan         nan         nan         nan         0.00020     0.00126     0.00488     nan
140     nan         nan         nan         nan         nan         0.00060     0.00397     nan
150     nan         nan         nan         nan         nan         0.00004     0.00310     nan
160     nan         nan         nan         nan         nan         nan         0.00230     nan
170     nan         nan         nan         nan         nan         nan         0.00157     nan
180     nan         nan         nan         nan         nan         nan         0.00084     nan

}\datatable

\begin{tikzpicture}

\begin{axis}[%
xlabel={$T$, $^\circ$C},
ylabel={Поверхневий натяг, Н/м},
width=0.75\linewidth,
height=0.5\linewidth,
scale only axis,
enlargelimits=false,
line join=round,
% === Налаштування сітки ===
grid = both,
grid style={line width=.1pt, draw=gray!10},
major grid style={line width=.2pt,draw=gray!50},
minor tick num = 5,
minor grid style = {line width=.1pt,draw=gray!10},
at={(7cm,0)},
xtick distance=20
]

\addplot[blue, thick, mark=none, smooth] table[x=x, y=y8] \datatable;
\addlegendentry{R134a \ce{CO2}}
\addplot[black, thick, mark=square, mark size=3pt, smooth] table[x=x, y=y1] \datatable;
\addlegendentry{R744}
\addplot[green!50!black, thick, mark=triangle, mark size=3pt, smooth] table[x=x, y=y2] \datatable;
\addlegendentry{R1234yf}
\addplot[violet, thick, mark=diamond, mark size=3pt, smooth] table[x=x, y=y3] \datatable;
\addlegendentry{R290 (Пропан)}
\addplot[blue!50!black, thick, mark=triangle, mark options={rotate=180}, mark size=3pt, smooth] table[x=x, y=y4] \datatable;
\addlegendentry{R1234ze}
\addplot[pink!80!black, thick, mark=o, mark size=3pt, smooth] table[x=x, y=y5] \datatable;
\addlegendentry{R600a (Ізобутан)}
\addplot[orange, thick, mark=*, mark size=3pt, smooth] table[x=x, y=y6] \datatable;
\addlegendentry{R600 (Бутан)}
\addplot[red, thick, mark=star, mark size=3pt, smooth] table[x=x, y=y7] \datatable;
\addlegendentry{R601 (Пентан)}


\end{axis}

\end{tikzpicture}%
\caption{Залежність коефіцієнту поверхневого натягу від температури для деяких холодоагентів --- речовин, що використовуються в холодильних установках.}
\label{pic:Pic9.5}
\end{figure}
%---------------------------------------------------------


%% --------------------------------------------------------
\section{Лапласів тиск}
%% --------------------------------------------------------

Незважаючи на малість, сили поверхневого натягу здатні створювати значні тиски. В цьому легко переконатися, якщо пригадати,  що всередині мильної бульбашки тиск більший, ніж ззовні. Так ситуація трапляється завжди, якщо існує викривлена поверхня між двома фазами або речовинами. Умова термодинамічної рівноваги передбачає рівність температур  систем, що знаходяться в контакті та їх хімічних потенціалів:

\begin{equation}
    T_1 = T_2, \quad \mu_1 = \mu_2.
\end{equation}

Тиски ж, як ми переконалися, можуть бути різні. Тоді при постійних $T$ та $V$ мінімум вільної енергії Гельмгольца відповідає умові:

\begin{equation}
    dF = 0 = -p_1dV_1 - p_2dV_2 + \sigma d\Sigma.
    \label{eq:min_free_energy}
\end{equation}

Вираз \eqref{eq:min_free_energy} означає, що будь-які варіації об’єму однієї із підсистем системи супроводжуються такими змінами поверхні, що в цілому $F$ лишається мінімальною. Оскільки $V = V_1 + V_2$, то $dV_1 = -dV_2$, тому різниця тисків в двох фазах (речовинах) дорівнює:

\begin{equation}
    \Delta p = p_1 - p_2 = \sigma \left( \frac{d\Sigma}{dV} \right).
    \label{eq:laplace_pressure}
\end{equation}


%---------------------------------------------------------
\begin{figure}[h!]\centering
\includegraphics[width=0.5\linewidth]{Pic9.6}
\caption{До виведення формули Лапласа.}
\label{pic:Pic9.6}
\end{figure}
%---------------------------------------------------------

Елементарний приріст об’єму складає $dV = R_I\theta_I\cdot R_{II}\theta_{II} dR$ (рис.~\ref{pic:Pic9.6}). При цьому приріст площі дорівнює:

\begin{equation*}
    d\Sigma =(R_I+dR)(R_{II}+dR)\theta_I\theta_{II}-R_IR_{II}\theta_I\theta_{II}=(R_I{+R}_{II})\theta_I\theta_{II}dR,
\end{equation*}
де $R_I$, $R_{II}$ --- локальні кривизни поверхні. Звідси:

\begin{equation*}
    \Delta p=\sigma\left(\frac{1}{R_I}+\frac{1}{R_{II}}\right)
\end{equation*}

Різниця тисків $\Delta p$, зумовлена кривизною поверхні, отримала назву \textit{Лапласового тиску}. Зокрема, для краплі сферичної форми радіусу $R$, що межує з повітрям:

\begin{equation}
    \Delta p = \frac{2\sigma}{R}.
    \label{eq:spherical_drop}
\end{equation}

%% --------------------------------------------------------
\section{Поверхневий натяг і явища змочування}
%% --------------------------------------------------------


Перш ніж перейти до розгляду капілярних явищ, розглянемо явище змочування. Вперше цю проблему поставили і дослідили Юнг та Лаплас в 1803-1804 роках. Розглянемо три фази що знаходяться в контакті один з одним. Наприклад це можуть бути три рідкі речовини, або два рідких і одно газоподібне (рис.~\ref{pic:Pic9.7a}). Умова механічної рівноваги між трьома фазами (або середовищами) полягає рівності нулю рівнодіючої сили, прикладеної до лінії трьох-фазного контакту (ЛТК), що проходить через точку О. В цю рівнодіючу входять сили поверхневого натягу, діючої вздовж границ і розділу  $\vec{\sigma}_{12}$, $\vec\sigma_{23}$, $\vec{\sigma}_{31}$, сили ваги рідини і сила гідростатичного тиску. Останні набагато менше дії сил поверхневого натягу,тому умови рівноваги для краплі малого розміру набувають вигляду:

Розглянемо покомпонентно векторне рівняння, отримаємо два складних рівняння:
\begin{equation}
\vec{\sigma}_{12} + \vec{\sigma}_{23} + \vec{\sigma}_{31} = 0
\end{equation}

Розписуючи покомпонентно векторне рівняння отримаємо два складних рівняння:
\begin{align}\label{eq:sigmas}
\sigma_{13} = \sigma_{12}\cos\theta_1 &+ \sigma_{23}\cos\theta_2 \quad (0 \leq \sigma_{12})^2 \leq \sigma_{23}^2 \\
\sigma_{12}\sin\theta_1 &= \sigma_{23}\sin\theta_2
\end{align}

Звідки отримаємо:
\begin{align*}
    |\cos\theta_1|=|(\sigma_{13}^2+\sigma_{12}^2-\sigma_{23}^2)/2\sigma_{13}\sigma_{12}|\le1\\
    |\cos\theta_2|=|(\sigma_{12}^2+\sigma_{23}^2-\sigma_{13}^2)/2\sigma_{12}\sigma_{23}| \le1.
\end{align*}

%=========================================================
\begin{figure}[h!]\centering
%---------------------------------------------------------
\begin{subfigure}[t]{0.45\linewidth}\centering
    \includegraphics[width=\linewidth]{Pic9.7a}
\caption{}
\label{pic:Pic9.7a}
\end{subfigure}
%---------------------------------------------------------
\begin{subfigure}[t]{0.45\linewidth}\centering
    \includegraphics[width=\linewidth]{Pic9.7b}
\caption{}
\label{pic:Pic9.7b}
\end{subfigure}
%---------------------------------------------------------
\caption{Умова механічної рівноваги між трьома середовищами: (\protect\subref{pic:Pic9.7a}) газоподібним (або рідким) $1$ та рідкими $2$ та $3$ і (\protect\subref{pic:Pic9.7b}) твердим (знизу) і рідким (крапля) та газоподібним або також рідким.}
\end{figure}
%=========================================================

Із того, що $|\cos\theta|<1$ слідує що рівновага можлива, якщо $\sigma_{13} < \sigma_{23}+\sigma_{12}$. В цьому випадку крапля поводиться подібно до краплі жиру в тарілці супу. В протилежному випадку рівновага між трьома середовищами неможлива і середовище $2$ займає проміжне положення між $1$ та $3$: на поверхні $1$ утворюється тонка плівка (рис.~\ref{pic:Pic9.8b}). В окремому випадку, коли середовище $3$ (рис.~\ref{pic:Pic9.7b}) є твердим, то вимірюючи кут змочування можна визначити коефіцієнти поверхневого натягу. Тоді із \eqref{eq:sigmas} достатньо тільки першого рівняння --- звідки знаходимо:
\begin{equation*}
    \cos\theta=\frac{\sigma_{13}-\sigma_{23}}{\sigma_{12}}
\end{equation*}
Якщо кут змочування дорівнює 180о (рис.~\ref{pic:Pic9.8a}, випадок $A$) говорять про повну незмочуваність (наприклад крапля ртуті на пластмасі), якщо нулю (рис.~\ref{pic:Pic9.8a}, випадок $D$), то про повну змочуваність (наприклад вода на знежиреному чистому склі). Проміжні випадки ($B$ та $C$ на рис.~\ref{pic:Pic9.8a}) --- часткова змочуваність. Якщо справа йдеться про змочуваність певної поверхні водою, то в першому випадку говорять про \textit{гідрофільність} поверхні, а в другому про \textit{гідрофобність}. Можливі звичайно і проміжні варіанти, як наприклад для води та тефлону  $\theta = 120^\circ$.  На практиці має місце гістерезис змочування. Якщо капля натікає на поверхню, то говорять про $\theta_\text{натікання}$, якщо витікає, то $\theta_\text{відтікання}$. Як правило $\theta_\text{натікання} > \theta_\text{відтікання}$ (рис.~\ref{pic:Pic9.8c}). Причини гістерезис різноманітні, до кінця незрозумілі, але однозначно залежать від шорсткості поверхні, її гетерогенності, механічних властивостей (деформаційний гістерезис), характеру взаємодії рідини із поверхнею (\textit{аутофобність}).

%=========================================================
\begin{figure}[h!]\centering
%---------------------------------------------------------
\begin{subfigure}[t]{0.32\linewidth}\centering
\includegraphics[height=0.5\linewidth]{Pic9.8a}
\caption{}
\label{pic:Pic9.8a}
\end{subfigure}
\quad%---------------------------------------------------------
\begin{subfigure}[t]{0.31\linewidth}\centering
\includegraphics[height=0.5\linewidth]{Pic9.8b}
\caption{}
\label{pic:Pic9.8b}
\end{subfigure}
%---------------------------------------------------------
\begin{subfigure}[t]{0.31\linewidth}\centering
\includegraphics[height=0.5\linewidth]{Pic9.8c}
\caption{}
\label{pic:Pic9.8c}
\end{subfigure}
%---------------------------------------------------------
\caption{Ілюстрація явищ змочування: (\protect\subref{pic:Pic9.8a}) --- різні випадки змочуванності;  (\protect\subref{pic:Pic9.8b}) --- неможливість співіснування трьох фаз (нафта на поверхні води);  (\protect\subref{pic:Pic9.8c}) --- кут натікання $2$ більше ніж кут відтікання $1$.}
\end{figure}
%=========================================================

%% --------------------------------------------------------
\section{Поверхневий натяг і капілярні явища}
%% --------------------------------------------------------


Як показує досвід більшість рідин зазнає підняття в капілярах (рис.~\ref{pic:Pic9.9}). Розрахуємося висоту підняття $h$ в капілярі радіуса $a$. Якщо капіляр достатньо малий, то меніск можна вважати сферичним. Тоді, якщо $\theta$ це кут змочування, то радіус кривизни дорівнює:

\begin{equation}
    R = \frac{a}{\cos\theta}.
    \label{eq:radius_curvature}
\end{equation}


%---------------------------------------------------------
\begin{figure}[h!]\centering
\includegraphics[width=0.5\linewidth]{Pic9.9}
\caption{Підняття рідини в капілярах. До виводу формули Журена.}
\label{pic:Pic9.9}
\end{figure}
%---------------------------------------------------------

Оскільки тиск в рідині $p_w$ менший, ніж ззовні на величину лапласового тиску,рідина в капілярах піднімається, тим вище, чим менше радіус капіляру. Надлишковий тиск компенсується силою тяжіння (вагою стовпа рідини). Тоді:

\begin{equation}
    \Delta p_0 - p_w = \rho g h,
    \label{eq:pressure_gravity}
\end{equation}
де $p_0$ --- атмосферний тиск, а $p_w$ --- тиск рідини на рівні меніска.

Порівнюючи \eqref{eq:laplace_pressure}) з \eqref{eq:pressure_gravity} і використовуючи \eqref{eq:radius_curvature}, знайдемо:

\begin{equation}
    h = \frac{2\sigma \cos\theta}{\rho g a}.
    \label{eq:jurin_law}
\end{equation}

Очевидно із формули, що при незмочуваності, рідина буде опускатися. Формула \eqref{eq:jurin_law} отримала назву \emph{формули Журена}.

\Closesolutionfile{answer}